\section{Metodología}

\subsection{Modelos evaluados}

Se seleccionaron cuatro modelos open-source, sin restricciones de licencia
para descarga (no gated), cubriendo un rango de tamaño de 0,36 a 1,71 mil
millones de parámetros (Tabla 1).

\begin{table}[htbp]
\centering
\caption{Modelos evaluados en este estudio.}
\label{tab:orig-modelos}
\begin{tabular}{lll}
\toprule
Modelo & Parámetros & Familia / origen \\
\midrule
SmolLM2-360M-Instruct & 0.36 B & HuggingFaceTB \\
Qwen2.5-0.5B-Instruct  & 0.49 B & Alibaba (Qwen) \\
Qwen2.5-1.5B-Instruct  & 1.54 B & Alibaba (Qwen) \\
SmolLM2-1.7B-Instruct  & 1.71 B & HuggingFaceTB \\
\bottomrule
\end{tabular}
\end{table}

Todos los modelos se ejecutaron con la librería transformers, en precisión
bfloat16 y sin cuantización adicional, sobre una máquina virtual con 2
núcleos de CPU y 7,8 GB de RAM, sin aceleración por GPU. La decodificación
fue determinista (greedy, sin muestreo, temperatura no aplicada) para
aislar la variable de tamaño del modelo del efecto de la aleatoriedad y
garantizar reproducibilidad.

\subsection{Dataset de comandos}

Se diseñó manualmente un conjunto de 32 comandos de domótica en español
rioplatense (voseo), organizado en seis categorías que capturan distintos
desafíos lingüísticos: encendido/apagado simple, ajuste con valor numérico
explícito (grados o porcentaje), ajuste relativo o cualitativo sin valor
numérico (``bajale un poco''), comandos sobre múltiples dispositivos
(``todas las luces''), consultas de estado, y comandos con contexto
conversacional o negación. Cada comando cuenta con una anotación de
referencia (ground truth) de cinco campos: intent (encender / apagar /
ajustar / consultar), dispositivo, ubicación, valor y unidad.

\subsection{Protocolo de prompting y esquema de salida}

Se utilizó un único prompt de sistema, idéntico para los cuatro modelos, que
especifica el esquema JSON de salida, el vocabulario de dispositivos y
ubicaciones esperado, y dos ejemplos few-shot (no incluidos en el conjunto
de evaluación). Se instruyó a los modelos a responder únicamente con el
objeto JSON, sin texto adicional. El uso de un prompt idéntico para todos
los modelos privilegia la comparabilidad por sobre el desempeño máximo de
cada modelo individual: es posible que un prompt afinado por modelo mejore
los resultados de cada uno, pero a costa de introducir una variable de
confusión en la comparación.

\subsection{Métricas}

Se reportan cinco métricas: (1) tasa de JSON válido (respuesta parseable
como objeto JSON); (2) coincidencia exacta (los cinco campos coinciden con
el ground truth); (3) exactitud por campo (intent, dispositivo, ubicación,
valor, unidad evaluados individualmente); (4) latencia promedio de
inferencia por comando, en segundos, en el entorno de CPU descripto; y (5)
exactitud por categoría lingüística del comando.

\subsection{Prompt utilizado}

Para favorecer la reproducibilidad del estudio, se reproduce a continuación
una versión condensada del prompt de sistema (idéntico en su versión
completa para los cuatro modelos), que en la implementación real incluía
además un segundo ejemplo few-shot omitido aquí por espacio:

\begin{quote}
Sos un módulo de interpretación de comandos de voz para un sistema domótico
en español rioplatense. Convertí el comando del usuario en UN objeto JSON
con estos campos: ``intent'' (encender/apagar/ajustar/consultar),
``dispositivo'' (luz, aire\_acondicionado, calefaccion, persiana, cortina,
tv, parlante, cerradura, ventilador, todos), ``ubicacion'' (living, cocina,
dormitorio, bano, garage, patio, entrada, todas), ``valor'' (número si el
comando lo da, si no null), ``unidad'' (``\%'' o ``$^{\circ}$C'' si
corresponde, si no null). Reglas: respondé SOLO el JSON, sin texto
adicional; si el comando ajusta sin dar un número, ``valor'' y ``unidad''
deben ser null. Ejemplo: Comando: ``Poné el aire de la cocina en 20
grados.'' \{"intent": "ajustar", "dispositivo": "aire\_acondicionado",
"ubicacion": "cocina", "valor": 20, "unidad": "$^{\circ}$C"\}
\end{quote}

Cada comando de evaluación se envió como un turno de usuario adicional
(Comando: ``...'') inmediatamente después de este prompt de sistema, sin
turnos previos de conversación.
